\scanpage{25}
RIGIDITÉ ET GROUPES ALG. \mspage{1}

\margin{``Theorem of the square''}
% slanted
\diagram{\begin{tikzcd}[ampersand replacement=\&,column sep=small,row sep=small]
{}\& X_1\times_S X_2\arrow[dl,"p_1"']\arrow[dr,"p_2"]\&{}\\
X_1\arrow[dr,"\varphi_1"',dash]\&{}\&X_2\arrow[dl,"\varphi_2",dash]\\
{}\&S\&{}
\end{tikzcd}}

\uline{Théorème} Soient $S$ un préschéma \inserted{\unsure{loc. noeth.}}, $\varphi_i:X_i\longrightarrow S$ $(i=1,2)$ deux préschémas de type fini sur $S$, \inserted{$\varepsilon_1$ une section de $X_1$ sur $S$} \deleted{$G$ un schéma en groupes loc. de type fini sur $S$, on suppose} soient $p_i$ \inserted{$(i=1,2)$} les deux projections de $X_1\times_S X_2$ sur les facteurs. On suppose a) $p_1$ est propre, et $p_{1*}(\underline{\mathrm O}_{X_1\times X_2})\xleftarrow{\sim}\underline{\mathrm O}_{X_1}$

\deleted{\unsure{b)} Pour tout voisinage ouvert $U$ de $\varepsilon_1(S)$, \unsure{$X_1$ est} \illegible{1} $U\times_S X_2$ est schématiquement dense dans $X_1\times_S X_2$}

\inserted{les fibres géom de $X_1/S$ sont connexes}
\deleted{\unsure{a)}} b) $p_1$ est universellement ouvert, \inserted{\uuline{ou} $X_1$ plat sur $S$, à fibres géom primaires}

[ \uline{Ex} $\varphi_2$ est propre, plat et $s\in S\Rightarrow k(s)\simeq H^0(X_{1s},\underline{\mathrm O}_{X_{1s}})$ ; \deleted{``\unsure{d'ailleurs|primaires}'' \unsure{les sous-algèbres art. intègres, etc|les sous-algèbres art. intègres|i.e. sans cycles ass. immergés}} $\varphi_1$ \deleted{est plat,} à fibres géom.$^{\mathrm t}$ irréductibles ]. Alors pour tt préschéma en groupes $G$ sur $S$, séparé sur $S$, et tt $S$-hom
\[
f:X_1\times_S X_2\longrightarrow G
\]
on a
\[
f=(f_1p_{\deleted{\unsure{2}}1})(\unsure{f_2|f}p_2)\tag*{(*)}
\]
où
\[
f_1:X_1\longrightarrow G,\qquad f_2:X_2\longrightarrow X
\]
sont des $S$-homomorphismes, déterminés de façon unique par la condition supplémentaire
\[
\deleted{f_2=f\circ(\varepsilon_1\times_S\operatorname{id}_{X_2})}
\]
\deleted{[ équivaut à :} $(**)$ \quad $f_1\varepsilon_1={}$ section unité, \deleted{pourvu que $\varphi_2$ soit un épim. ]} \inserted{et alors $(***)\ f_2=f\circ(\varepsilon_1\times\operatorname{id}_{X_2})$}

\margin{Attention ! on ne suppose pas la loi de composition commutative}
% slanted
\margin{plus haut}
% slanted
\uline{Corollaire} Supposons de plus que $X_2$ soit muni d'une section $\varepsilon_2$ sur $S$. Alors $f$ se met de façon unique sous la forme
\[
f_1p_1+f_2p_2+\sigma\varphi
\]
\inserted{tels que $f_i\varepsilon_i={}$ section unité;} où $f_i:X_i\longrightarrow G$ sont des $S$-morphismes, $\sigma:S\longrightarrow G$ une section de $G$ sur $S$, [ $\varphi$ désignant \deleted{$\varphi_1\times_S\varphi_2$} la projection $X_1\times_S X_2\longrightarrow S$ ].
% archive sheet number: 24
